\subsection*{Fuentes de datos}
Se utilizaron los archivos del dataset Olist: \texttt{orders}, \texttt{customers}, \texttt{payments}, \texttt{reviews} e \texttt{items}. Para evitar duplicación por relaciones 1:N, se agregaron pagos, reviews e items por \texttt{order\_id} antes de construir el dataset a nivel cliente.

\subsection*{Construcción del label (churn)}
Se definió una fecha de \textit{snapshot} y una ventana futura $W=90$ días. El churn se define como:
\[
churn=1 \Leftrightarrow \text{el cliente no realiza compras en } (snapshot, snapshot+W]
\]
y $churn=0$ en caso contrario (recompra). Para evitar sesgos por clientes históricamente inactivos, se filtra la población a clientes con compra reciente dentro de un \textit{lookback} previo al snapshot.

\subsection*{Evidencia de integridad del dataset}
En el preprocesamiento a nivel orden se verificó que no hubiera duplicación de \texttt{order\_id} (duplicados = 0), evitando inflación artificial de montos o conteos.